% ---- CLASSE E PACOTES ESSENCIAIS --------------
\documentclass[12pt, a4paper, twoside, openright]{report}


% Encoding e língua
\usepackage[utf8]{inputenc}
\usepackage[T1]{fontenc}
\usepackage[portuguese]{babel}

% Geometria da página — margens do template Word
\usepackage[
  a4paper,
  top=2cm,
  bottom=2cm,
  left=3.0cm,
  right=2.5cm,
  headsep=0.75cm,
  footskip=1.0cm
]{geometry}

% Fontes
\usepackage[scaled=0.95]{helvet}
\usepackage{mathptmx}
\renewcommand{\familydefault}{\rmdefault}
\usepackage{microtype}

% Cores
\usepackage[dvipsnames,svgnames,x11names]{xcolor}
\definecolor{isepdark}{HTML}{9A1A24}
\definecolor{isepgray}{HTML}{666666}

% Espaçamento
\usepackage{setspace}
\setstretch{1.5}

% Parágrafo
\usepackage{parskip}
\setlength{\parskip}{6pt}
\setlength{\parindent}{0pt}

% Cabeçalho e rodapé
\usepackage{fancyhdr}
\usepackage{emptypage}

% Títulos
\usepackage{titlesec}

% Índices
\usepackage{tocloft}

% Figuras e tabelas
\usepackage{graphicx}
\usepackage{float}
\usepackage{caption}
\usepackage{subcaption}
\usepackage{booktabs}
\usepackage{longtable}
\usepackage{array}
\usepackage{multirow}
\usepackage{tabularx}

% Código fonte
\usepackage{listings}
\usepackage{courier}

% Hiperligações
\usepackage[
  colorlinks=true,
  linkcolor=isepdark,
  citecolor=isepdark,
  urlcolor=isepdark,
  bookmarks=true,
  bookmarksnumbered=true
]{hyperref}

% Glossário
\usepackage[acronym, toc, nonumberlist]{glossaries}

% Matemática
\usepackage{amsmath, amssymb}
\usepackage{siunitx}

% Miscelânea
\usepackage{csquotes}
\usepackage{enumitem}
\usepackage{pdfpages}

% Todonotes — para marcar os placeholders de imagens
\usepackage[disable]{todonotes} % Mudar para [disable] na versão final para esconder os todos

% =============================================================================
% FORMATAÇÃO DE TÍTULOS
% =============================================================================
\titleformat{\chapter}[block]
  {\sffamily\bfseries\fontsize{20}{24}\selectfont\color{isepdark}}
  {\thechapter\quad}
  {0pt}
  {}
\titlespacing*{\chapter}{0pt}{2.0cm}{0.5cm}

\titleformat{\section}[block]
  {\sffamily\bfseries\fontsize{14}{17}\selectfont\color{isepdark}}
  {\thesection\quad}
  {0pt}
  {}
\titlespacing*{\section}{1.02cm}{0.5cm}{0.2cm}

\titleformat{\subsection}[block]
  {\sffamily\bfseries\fontsize{12}{15}\selectfont\color{isepdark}}
  {\thesubsection\quad}
  {0pt}
  {}
\titlespacing*{\subsection}{1.02cm}{0.4cm}{0.1cm}

\titleformat{\subsubsection}[block]
  {\sffamily\bfseries\fontsize{11}{14}\selectfont\color{isepdark}}
  {\thesubsubsection\quad}
  {0pt}
  {}
\titlespacing*{\subsubsection}{1.02cm}{0.3cm}{0.1cm}

% =============================================================================
% FORMATAÇÃO DO ÍNDICE
% =============================================================================
\setcounter{tocdepth}{3}
\setcounter{secnumdepth}{3}
\renewcommand{\cftchapfont}{\sffamily\bfseries}
\renewcommand{\cftsecfont}{\sffamily}
\renewcommand{\cftsubsecfont}{\sffamily}
\renewcommand{\cftsubsubsecfont}{\sffamily}

% =============================================================================
% CABEÇALHO E RODAPÉ
% =============================================================================
\newcommand{\tituloprojeto}{SMARTER HUB}
\newcommand{\nomeautor}{Patrick Dias Mendes da Costa}

\pagestyle{fancy}
\fancyhf{}
\fancyhead[LO]{\small\sffamily\color{isepgray}\leftmark}
\fancyhead[RO]{\small\sffamily\color{isepgray}\tituloprojeto}
\fancyhead[LE]{\small\sffamily\color{isepgray}\tituloprojeto}
\fancyhead[RE]{\small\sffamily\color{isepgray}\leftmark}
\fancyfoot[LO,LE]{\small\sffamily\color{isepgray}\nomeautor}
\fancyfoot[RO,RE]{\small\sffamily\color{isepgray}\thepage}
\renewcommand{\headrulewidth}{0.4pt}
\renewcommand{\footrulewidth}{0pt}

\fancypagestyle{plain}{
  \fancyhf{}
  \fancyfoot[RO,RE]{\small\sffamily\color{isepgray}\thepage}
  \renewcommand{\headrulewidth}{0pt}
}

\renewcommand{\chaptermark}[1]{\markboth{#1}{}}

% =============================================================================
% LEGENDAS
% =============================================================================
\captionsetup[figure]{
  labelfont={bf,sf}, textfont=sf, labelsep=endash,
  position=below, justification=centering
}
\captionsetup[table]{
  labelfont={bf,sf}, textfont=sf, labelsep=endash,
  position=above, justification=centering
}

\addto\captionsportuguese{
  \renewcommand{\figurename}{Figura}
  \renewcommand{\tablename}{Tabela}
  \renewcommand{\contentsname}{Índice}
  \renewcommand{\listfigurename}{Índice de Figuras}
  \renewcommand{\listtablename}{Índice de Tabelas}
  \renewcommand{\bibname}{Bibliografia}
  \renewcommand{\appendixname}{Anexo}
  \renewcommand{\chaptername}{}
}

% =============================================================================
% ESTILO CÓDIGO FONTE
% =============================================================================
\lstset{
  basicstyle=\ttfamily\small,
  breaklines=true,
  breakatwhitespace=true,
  frame=single,
  framerule=0.4pt,
  rulecolor=\color{isepgray!40},
  numbers=left,
  numberstyle=\tiny\color{isepgray},
  numbersep=8pt,
  xleftmargin=14pt,
  xrightmargin=4pt,
  keywordstyle=\color{isepdark}\bfseries,
  stringstyle=\color{teal},
  commentstyle=\color{isepgray}\itshape,
  showstringspaces=false,
  language=[Sharp]C,
}

% =============================================================================
% NUMERAÇÃO DE PÁGINAS
% =============================================================================
\newcommand{\frontmatter}{\pagenumbering{roman}}
\newcommand{\mainmatter}{\pagenumbering{arabic}}

% =============================================================================
% ANEXOS
% =============================================================================
\usepackage[title,titletoc]{appendix}
\renewcommand{\appendixname}{Anexo}

% =============================================================================
% DADOS DO RELATÓRIO
% =============================================================================
\newcommand{\licenciatura}{Licenciatura em Engenharia de Sistemas}
\newcommand{\titulorelatorio}{SMARTER HUB: Plataforma Integrada de Gestão de Colaboradores}
\newcommand{\nomeorganizacao}{Tlantic Portugal - Sistemas de Informação, S.A.}
\newcommand{\dataapresentacao}{junho de 2026}
\newcommand{\orientadorisep}{Marílio Cardoso}
\newcommand{\supervisorexterno}{Camila Teixeira}

\renewcommand{\tituloprojeto}{\titulorelatorio}
\renewcommand{\nomeautor}{Patrick Dias Mendes da Costa}

% =============================================================================
% GLOSSÁRIO
% =============================================================================
\makeglossaries

\newacronym{proes}{PROES}{Unidade Curricular de Projeto/Estágio}
\newacronym{si}{SI}{Sistema Internacional de Unidades}
\newacronym{uml}{UML}{Linguagem de Modelagem Universal}
\newacronym{api}{API}{Application Programming Interface}
\newacronym{rest}{REST}{Representational State Transfer}
\newacronym{jwt}{JWT}{JSON Web Token}
\newacronym{orm}{ORM}{Object-Relational Mapper}
\newacronym{crud}{CRUD}{Create, Read, Update, Delete}
\newacronym{rh}{RH}{Recursos Humanos}
\newacronym{e2e}{E2E}{End-to-End}
\newacronym{dto}{DTO}{Data Transfer Object}
\newacronym{iban}{IBAN}{International Bank Account Number}
\newacronym{nif}{NIF}{Número de Identificação Fiscal}
\newacronym{niss}{NISS}{Número de Identificação da Segurança Social}
\newacronym{ptbr}{PT/BR}{Portugal e Brasil}
\newacronym{mvp}{MVP}{Minimum Viable Product}
\newacronym{moscow}{MoSCoW}{Must have, Should have, Could have, Won't have}
\newacronym{cpf}{CPF}{Cadastro de Pessoas Físicas}
\newacronym{pis}{PIS}{Programa de Integração Social}

% =============================================================================
% DOCUMENTO
% =============================================================================
\begin{document}

% -------------------------
% CAPA
% -------------------------
\begin{titlepage}
  \thispagestyle{empty}
  \centering

  \vspace*{1cm}
  {\large\sffamily\bfseries\color{isepdark} \licenciatura}

  \vspace{3cm}
  {\LARGE\sffamily\bfseries \titulorelatorio}

  \vspace{1.5cm}
  {\large\sffamily \nomeorganizacao}

  \vfill

  {\large\sffamily \nomeautor}

  \vspace{1cm}
  {\sffamily \dataapresentacao}

  \vspace{1cm}
  {\sffamily Orientador ISEP: \orientadorisep}\\
  {\sffamily Supervisor Externo: \supervisorexterno}

  \vspace{1cm}
\end{titlepage}

% -------------------------
% PRÉ-MATÉRIA
% -------------------------
\frontmatter
\pagestyle{fancy}

% - Agradecimentos -
\chapter*{Agradecimentos}
\addcontentsline{toc}{chapter}{Agradecimentos}
\fancyhead[LO,RE]{Agradecimentos}

Agradeço à Tlantic Portugal pela oportunidade de realizar o estágio em contexto empresarial
real, com exposição a problemas concretos de engenharia de software e à operação de
Recursos Humanos numa empresa em crescimento. A confiança depositada num estudante
para construir, de raiz, uma plataforma que vai ser efetivamente utilizada é algo que não é
garantido em muitos estágios curriculares, e que tornou esta experiência genuinamente
desafiante.

Agradeço ao orientador, Marílio Cardoso, pelo acompanhamento académico, pela validação
metodológica e pelo apoio na estruturação técnica e documental do trabalho.

Agradeço à supervisora externa, Camila Teixeira, pela disponibilidade permanente,
pelo enquadramento funcional dos requisitos de negócio e pela paciência em iterações
que por vezes exigiram rever decisões que pareciam estar fechadas. O seu olhar sobre o
produto foi determinante para que as escolhas técnicas ficassem alinhadas com o uso real.

Agradeço ainda à equipa com quem colaborei diariamente, pela partilha de conhecimento
e pela disponibilidade em validar fluxos e comportamentos da plataforma.

\vspace{1cm}

% - Resumo -
\chapter*{Resumo}
\addcontentsline{toc}{chapter}{Resumo}
\fancyhead[LO,RE]{Resumo}

Este relatório descreve o trabalho desenvolvido no estágio curricular realizado na Tlantic
Portugal, no âmbito da unidade curricular PROES do ano letivo 2025/2026. O projeto
consistiu no desenvolvimento do SMARTER HUB, uma plataforma \textit{web} integrada para
gestão operacional de pessoas, com foco em perfis de colaboradores, férias e ausências,
e governação de acessos.

O contexto de partida era marcado por fragmentação operacional: processos de aprovação
de férias distribuídos por email, folhas de cálculo sem controlo de versão e validações
informais que tornavam difícil responder à pergunta mais simples que um gestor de equipa
pode colocar - \textit{``quem está de férias esta semana?''} A ausência de rastreabilidade
e a diferença de legislação laboral entre Portugal e Brasil, países onde a Tlantic opera,
acrescentavam camadas adicionais de complexidade a um problema que existia, em grande
medida, por falta de um sistema que o centralizasse.

O projeto foi desenvolvido como trabalho individual ao longo de três meses, adotando
uma metodologia iterativa com reuniões semanais de acompanhamento. A solução
implementada é uma aplicação \textit{full-stack}: React com TypeScript no \textit{frontend},
Node.js com Express no \textit{backend}, e PostgreSQL com Prisma ORM na camada de
persistência. A plataforma suporta fluxos de aprovação com hierarquias distintas para
Portugal e Brasil, um modelo de permissões com trilho auditável de concessões e
revogações, e mecanismos de versionamento de pedidos que garantem rastreabilidade
completa do histórico de cada decisão.

Como principais resultados, destaca-se a entrega de quatro módulos funcionais em
ambiente de pré-produção, com validação técnica por testes unitários, de integração e
\textit{end-to-end}. O trabalho consolidou competências de desenvolvimento \textit{full-stack}
em ambiente profissional real, com impacto direto na qualidade do produto e no
alinhamento entre engenharia de software e necessidades operacionais de negócio.

\bigskip
\noindent\textbf{Palavras-chave:}\quad gestão de pessoas; férias e ausências; governação de
acessos; rastreabilidade; engenharia de software.

\newpage

% - Índice geral -
\fancyhead[LO,RE]{{\color{isepdark}Índice}}
\tableofcontents

\newpage

% - Índice de Figuras -
\fancyhead[LO,RE]{Índice de Figuras}
\listoffigures
\addcontentsline{toc}{chapter}{Índice de Figuras}

\newpage

% - Índice de Tabelas -
\fancyhead[LO,RE]{Índice de Tabelas}
\listoftables
\addcontentsline{toc}{chapter}{Índice de Tabelas}

\newpage

% - Notação e Glossário -
\chapter*{Notação e Glossário}
\addcontentsline{toc}{chapter}{Notação e Glossário}
\fancyhead[LO,RE]{Notação e Glossário}

\textit{Esta secção apresenta os acrónimos e conceitos técnicos utilizados no corpo do
texto, ordenados alfabeticamente.}


\begin{table}[!h]
\begin{tabular}{lp{0.84\textwidth}}
  \toprule
  \textbf{Termo} & \textbf{Definição} \\
  \midrule
  \textbf{API}        & \textit{Application Programming Interface} - conjunto de endpoints REST que expõe a lógica de negócio do \textit{backend} ao \textit{frontend}. \\[4pt]
  \textbf{CPF}        & \textit{Cadastro de Pessoas Físicas} - número de identificação fiscal brasileiro, equivalente ao NIF português. \\[4pt]
  \textbf{CRUD}       & \textit{Create, Read, Update, Delete} - operações básicas de persistência de dados. \\[4pt]
  \textbf{DTO}        & \textit{Data Transfer Object} - estrutura de dados usada para validar e transportar informação entre camadas da aplicação. \\[4pt]
  \textbf{E2E}        & \textit{End-to-End} - testes que simulam o comportamento real do utilizador através de toda a \textit{stack}. \\[4pt]
  \textbf{IBAN}       & \textit{International Bank Account Number} - número de conta bancária internacional. \\[4pt]
  \textbf{JWT}        & \textit{JSON Web Token} - mecanismo de autenticação \textit{stateless} utilizado para gerir sessões na API. \\[4pt]
  \textbf{MoSCoW}     & Método de priorização de requisitos: \textit{Must have, Should have, Could have, Won't have}. \\[4pt]
  \textbf{NIF}        & Número de Identificação Fiscal (Portugal). \\[4pt]
  \textbf{NISS}       & Número de Identificação da Segurança Social (Portugal). \\[4pt]
  \textbf{ORM}        & \textit{Object-Relational Mapper} - camada de abstração sobre a base de dados; neste projeto, o Prisma. \\[4pt]
  \textbf{PIS}        & \textit{Programa de Integração Social} - número de identificação trabalhista brasileiro. \\[4pt]
  \textbf{PROES}      & Unidade Curricular de Projeto/Estágio. \\[4pt]
  \textbf{PT/BR}      & Abreviatura usada no projeto para referir as especificidades dos contextos legislativos de Portugal e do Brasil. \\[4pt]
  \textbf{REST}       & \textit{Representational State Transfer} - estilo arquitetural para APIs HTTP. \\[4pt]
  \textbf{RH}         & Recursos Humanos. \\[4pt]
  \textbf{UML}        & \textit{Unified Modeling Language} - linguagem de modelação utilizada em diagramas de casos de uso e de domínio. \\[4pt]
  \bottomrule
\end{tabular}
\end{table}


% -------------------------
% CORPO
% -------------------------
\mainmatter
\pagestyle{fancy}

% =========================================================================
\chapter{Introdução}
\label{cap:introducao}
% =========================================================================

O presente relatório documenta o trabalho desenvolvido durante o estágio curricular realizado
na Tlantic Portugal - Sistemas de Informação, S.A., entre abril e junho de 2026. O projeto
consistiu no desenvolvimento do \textbf{SMARTER HUB}, uma plataforma \textit{web} para
centralizar e digitalizar os processos operacionais de gestão de pessoas da empresa.

A escolha do tema surgiu de uma necessidade real: antes do SMARTER HUB, a gestão de
férias, perfis de colaboradores e permissões de acesso era feita de forma dispersa, sem
um sistema único que desse visibilidade aos gestores e garantisse rastreabilidade das
decisões. O desafio era construir algo que resolvesse esse problema num contexto
empresarial real - com utilizadores reais, regras de negócio complexas e dois contextos
legislativos diferentes.

\section{Enquadramento}

A Tlantic Portugal opera num mercado onde a eficiência operacional é determinante.
Com uma equipa distribuída por Portugal e Brasil, a gestão de pessoas envolve não apenas
a diversidade cultural, mas também diferenças legislativas significativas entre os dois países.
Antes do SMARTER HUB, esta gestão dependia de folhas de cálculo partilhadas, pedidos
por email e validações informais entre gestores e equipas de RH.

O projeto foi enquadrado como uma iniciativa de centralização e digitalização de fluxos
críticos de RH, com três eixos principais: uniformização de regras, melhoria da
rastreabilidade das decisões e redução da carga administrativa.

\section{Apresentação do projeto}

O SMARTER HUB é uma aplicação \textit{web} desenvolvida de raiz como projeto de estágio,
destinada a suportar a gestão integrada de pessoas na Tlantic. A plataforma cobre:

\begin{itemize}
  \item gestão de perfis de colaboradores com campos específicos por país (PT/BR);
  \item fluxos de férias e ausências com estados explícitos, versionamento e deteção de
        sobreposições;
  \item governação de permissões de acesso com trilho auditável;
  \item sistema de notificações para eventos críticos;
  \item mecanismos de exportação para suporte à decisão.
\end{itemize}

O desenvolvimento foi realizado de forma individual, com acompanhamento semanal da
supervisora externa. No final do período de estágio, a plataforma encontrava-se em
fase de pré-produção, com entrega prevista a uma equipa piloto para testes.

\section{Definição dos Objetivos}
\label{sec:objetivos}

\textbf{Objetivos gerais}

\begin{itemize}
  \item Desenvolver uma plataforma integrada para os processos de gestão de pessoas
        no contexto da Tlantic.
  \item Reduzir a dependência de processos manuais e ferramentas dispersas.
  \item Garantir rastreabilidade e governança em operações sensíveis de RH.
\end{itemize}

\textbf{Objetivos específicos}

\begin{itemize}
  \item Centralizar a gestão de informação de colaboradores numa única aplicação \textit{web}.
  \item Implementar fluxos de férias/ausências com estados explícitos, validações de
        capacidade de equipa e histórico de pedidos.
  \item Suportar dois contextos legislativos distintos (Portugal e Brasil) no mesmo sistema.
  \item Estruturar um modelo de permissões com rastreabilidade de concessões e revogações.
  \item Disponibilizar mecanismos de reporte e exportação de dados.
  \item Garantir qualidade técnica através de testes unitários, de integração e \textit{end-to-end}.
  \item Criar uma base evolutiva que suporte novos requisitos sem reestruturação completa.
\end{itemize}

\section{Calendarização e Metodologia}

O estágio decorreu entre 6 de abril e finais de junho de 2026, ao longo de doze semanas.
O Gantt \textit{chart} na Figura~\ref{fig:gantt} ilustra a distribuição temporal das principais
fases do projeto.

\begin{figure}[H]
  \centering
  % INSTRUÇÃO: Substituir pela exportação do Gantt chart real (PNG ou PDF)
  % \includegraphics[width=\textwidth]{imagens/gantt_chart.png}
  \missingfigure[figwidth=\textwidth]{Gantt chart com o plano de trabalhos de PROES}
  \caption{Plano de trabalhos do estágio PROES 2025/2026}
  \label{fig:gantt}
\end{figure}

A metodologia adotada foi \textbf{iterativa e incremental}. Cada semana terminava com
uma reunião de acompanhamento com a supervisora Camila Teixeira, onde se avaliava
o trabalho realizado, reajustavam-se prioridades e clarificavam-se requisitos para a
semana seguinte. Esta cadência permitiu responder a novos requisitos que surgiram 
ao longo do projeto, nomeadamente funcionalidades que só ficaram claras após a
implementação de outras, sem comprometer a estabilidade do que já estava entregue.

Na prática, cada incremento seguiu três etapas:

\begin{enumerate}
  \item \textbf{Levantamento e especificação} - clarificação de requisitos com a supervisora
        e análise das regras de negócio envolvidas;
  \item \textbf{Implementação e teste} - desenvolvimento da funcionalidade, cobertura
        com testes e revisão do código;
  \item \textbf{Validação e adaptação} - demonstração do resultado e incorporação
        de \textit{feedback}.
\end{enumerate}

\section{Motivação e Contributos}

Do ponto de vista organizacional, o principal contributo foi a disponibilização de uma
plataforma que substitui um conjunto de processos manuais por fluxos digitais,
auditáveis e com regras de negócio explícitas. A redução da carga administrativa
em tarefas repetitivas, como a produção do mapa anual de férias, e a centralização
da informação de colaboradores são os resultados com maior impacto imediato.

Do ponto de vista individual, este estágio representou a primeira experiência de
desenvolvimento de software em contexto profissional real. A necessidade de construir
algo que será efetivamente utilizado por utilizadores reais, dados sensíveis
e regras de negócio que não podem ser aproximadas foi o principal motor de
aprendizagem, e também a principal fonte de dificuldade.

\section{Utilização de Conteúdo Gerado por IA}

No âmbito do projeto foi utilizado o GitHub Copilot como ferramenta de apoio à
produtividade, principalmente na implementação de código, na aceleração de tarefas
de refatoração e na criação de testes. A utilização foi sempre supervisionada:
as sugestões geradas foram avaliadas criticamente, adaptadas ao contexto específico
do SMARTER HUB e validadas antes de serem integradas.

É importante distinguir entre apoio instrumental e autoria intelectual. As decisões
de arquitetura, a modelação do domínio, as regras de negócio PT/BR e o texto académico
são da responsabilidade do autor. O Copilot funcionou como um acelerador de execução,
não como substituto de raciocínio.

\begin{table}[H]
  \caption{Utilização de IA no projeto}
  \label{tab:ia}
  \small\sffamily
  \begin{tabularx}{\textwidth}{|l|l|X|l|}
    \hline
    \textbf{Categoria} & \textbf{Ferramenta} & \textbf{Justificação e salvaguarda} & \textbf{Capítulos} \\
    \hline
    Texto      & GitHub Copilot & Apoio à estruturação de texto técnico; revisão e validação integral pelo autor. & 1, 2 \\
    \hline
    Imagem     & N/A & Sem geração automática de imagens; todos os diagramas foram criados pelo autor. & - \\
    \hline
    Código     & GitHub Copilot & Sugestões de implementação, testes e refatoração; todas revisadas e adaptadas. & 3, 4 \\
    \hline
    Vídeo      & N/A & Não aplicável. & - \\
    \hline
    Outros     & N/A & Não aplicável. & - \\
    \hline
  \end{tabularx}
\end{table}

\section{Organização do Relatório}

O relatório está organizado em cinco capítulos. O Capítulo~\ref{cap:introducao} apresenta
o enquadramento, objetivos e metodologia. O Capítulo~\ref{cap:contexto} descreve o
contexto organizacional, o problema identificado e a decisão tecnológica. O
Capítulo~\ref{cap:atividades} detalha as atividades desenvolvidas ao longo do estágio.
O Capítulo~\ref{cap:arquitetura} consolida a arquitetura da solução e a estratégia de
validação técnica. Por fim, o Capítulo~\ref{cap:conclusao} reúne as conclusões, os
objetivos realizados e as direções de trabalho futuro.

% =========================================================================
\chapter{Contexto}
\label{cap:contexto}
% =========================================================================

\section{Organização e Área de Negócio}

A Tlantic Portugal - Sistemas de Informação, S.A. é uma empresa portuguesa de
desenvolvimento de \textit{software} fundada no Brasil em 2004, com presença consolidada
em Portugal desde 2007. Com cerca de duzentos colaboradores e clientes como a Sonae,
Auchan, Lidl, o Grupo Pão de Açúcar e a Parfois, a empresa atua como parceira
tecnológica de grandes organizações nos setores do retalho, indústria e saúde.

O posicionamento da Tlantic centra-se na eficiência operacional: os seus produtos ajudam
os clientes a gerir pessoas, operações e ponto de venda de forma integrada. A empresa
opera com equipas em Portugal e no Brasil, o que cria uma dualidade legislativa laboral
que está diretamente presente no contexto do SMARTER HUB. A mesma plataforma
tem de suportar regras de férias, aprovações e contratos diferentes, consoante o país do
colaborador.

\section{Problema Identificado}

Antes do SMARTER HUB, a gestão de pessoas na Tlantic dependia de um conjunto
heterogéneo de ferramentas sem integração entre si. Os pedidos de férias eram enviados
através de uma plataforma, aprovados informalmente e registados em folhas de cálculo partilhadas. As
permissões de acesso a sistemas eram atribuídas sem registo formal de justificativas.
Os dados de colaboradores estavam distribuídos por diferentes fontes, sem um modelo
único de atualização.

As principais limitações observadas foram:

\begin{itemize}
  \item aprovações de férias sem rastreabilidade: não era possível saber, de forma
        simples, quem aprovou o quê, quando e com que informação;
  \item ausência de validação automática de sobreposições de ausências dentro da mesma
        equipa;
  \item dificuldade em aplicar políticas laborais distintas para colaboradores portugueses
        e brasileiros no mesmo fluxo;
  \item dados de perfil desatualizados por falta de um mecanismo central de gestão;
  \item permissões de acesso sem histórico de concessão ou revogação.
\end{itemize}

\section{Análise de Causas}

A análise inicial do problema identificou quatro causas estruturais:

\begin{description}
  \item[Dados dispersos.] A informação de colaboradores estava fragmentada por email,
        ficheiros partilhados e sistemas de terceiros, sem um modelo único de atualização
        ou validação.
  \item[Ausência de processo formal.] Os fluxos de aprovação de férias e de atribuição
        de permissões funcionavam por convenção informal, sem estados definidos nem
        histórico de decisões.
  \item[Falta de validação automática.] Não existia mecanismo que detetasse, por exemplo,
        que dois colaboradores da mesma equipa tinham férias sobrepostas num período crítico.
  \item[Dualidade legislativa não suportada.] As diferenças entre a legislação laboral
        portuguesa e brasileira - em matéria de elegibilidade, prazos e níveis de aprovação
        - eram geridas ad hoc, sem regras explícitas no sistema.
\end{description}

\section{Perspetiva do Utilizador}

O perfil de utilizador com maior atrito operacional identificado foi o \textbf{gestor de equipa}.
Na prática, um gestor precisa responder rapidamente a pedidos de férias sem comprometer
a continuidade operacional da equipa. O problema não era apenas a demora no processo,
mas a incerteza: ao aprovar um pedido, o gestor não tinha acesso imediato à informação
de quem mais estava de férias naquele período, o que tornava cada decisão potencialmente
arriscada.

Do lado da equipa de RH, uma das grandes queixas era a produção do mapa anual de férias, uma tarefa que exigia consolidar manualmente informação de múltiplas fontes.

\section{Estado da Arte e Opção Tecnológica}

Antes de iniciar o desenvolvimento, foi feita uma análise comparativa entre adquirir uma
solução de mercado (\textit{buy}) e desenvolver uma solução orientada ao contexto (\textit{build}).
As soluções HRIS (\textit{Human Resource Information Systems}) de mercado, como o Personio,
o BambooHR ou o HiBob, oferecem funcionalidades de gestão de pessoas abrangentes,
mas a sua parametrização para suportar simultaneamente a legislação laboral portuguesa
e brasileira em um único sistema, com as especificidades de hierarquia e aprovação da
Tlantic, exigiria customizações extensas e dependência de um fornecedor externo.

A opção pelo desenvolvimento interno foi justificada por três razões principais:

\begin{itemize}
  \item \textbf{Controlo sobre regras PT/BR.} A coexistência de dois contextos legislativos
        com lógicas de aprovação distintas exige implementação customizada que não
        está disponível \textit{out-of-the-box} nas soluções de mercado analisadas.
  \item \textbf{Controlo sobre o \textit{roadmap}.} O desenvolvimento interno permite evoluir
        a plataforma em função das necessidades reais da Tlantic, sem dependência de
        ciclos de atualização de fornecedores.
  \item \textbf{Alinhamento com a estratégia tecnológica.} A Tlantic opera maioritariamente
        em PHP nas suas plataformas internas, mas tem a intenção de evoluir para tecnologias
        mais modernas. O SMARTER HUB foi desenhado intencionalmente com uma
        \textit{stack} diferente: React, TypeScript, Node.js e PostgreSQL, como
        exercício prático nessa direção.
\end{itemize}

\section{Stakeholders Principais}

Os atores com maior influência no projeto foram:

\begin{itemize}
  \item \textbf{Supervisora externa (Camila Teixeira)} - interlocutora principal para
        requisitos de negócio, validação funcional e priorização do \textit{backlog};
  \item \textbf{Equipa de RH} - utilizadores finais dos módulos de férias, perfis
        e permissões; principal fonte de requisitos funcionais;
  \item \textbf{Gestores de equipa} - aprovadores nos fluxos de férias; utilizadores
        do módulo de visibilidade de ausências;
  \item \textbf{Colaboradores} - utilizadores finais do perfil próprio e do fluxo
        de submissão de pedidos;
  \item \textbf{Equipa técnica} - responsável pela revisão e integração futura
        da plataforma no ambiente de produção da Tlantic.
\end{itemize}

% =========================================================================
\chapter{Atividades Desenvolvidas}
\label{cap:atividades}
% =========================================================================

Este capítulo descreve o trabalho realizado ao longo do estágio: desde o levantamento
de requisitos até à entrega dos módulos funcionais, passando pelas decisões técnicas
que moldaram a implementação. A apresentação segue a ordem cronológica do projeto,
com ênfase nos aspetos que foram mais relevantes do ponto de vista técnico e de
aprendizagem.

\section{Planeamento e Critérios de Aceitação}

O desenvolvimento foi organizado em \textit{sprints} semanais, cada um com objetivos
definidos no início da semana e revistos na reunião de fim de semana com a supervisora.
Esta cadência permitiu manter o alinhamento entre o que era tecnicamente implementado
e o que era funcionalmente esperado, e foi fundamental para gerir a complexidade
crescente dos requisitos.

Cada funcionalidade foi considerada entregue quando satisfazia os seguintes critérios:

\begin{itemize}
  \item os requisitos funcionais definidos para o fluxo estavam implementados
        e comportavam-se conforme o esperado;
  \item as regras de negócio e restrições de integridade de dados estavam validadas;
  \item a interface era consistente com os padrões visuais do restante SMARTER HUB;
  \item existia cobertura de testes para os fluxos críticos da funcionalidade;
  \item a supervisora ou outro utilizador-chave reconhecia o comportamento implementado
        como correto.
\end{itemize}

\section{Análise e Especificação de Requisitos}

A fase de análise decorreu nas primeiras semanas de estágio e continuou de forma
intermitente ao longo de todo o projeto, à medida que novos requisitos foram emergindo.
O processo de levantamento combinou reuniões com a supervisora, análise dos processos
informais existentes e, em alguns casos, discussão direta com utilizadores finais para
perceber onde estava o atrito.

\subsection{Inventário Funcional}

O inventário de requisitos foi organizado em quatro domínios principais:

\begin{description}
  \item[Perfil de colaborador.] Dados pessoais, informação contratual, regime horário,
        \gls{iban}, \gls{nif}/\gls{niss} para Portugal e \gls{cpf}/\gls{pis} para o
        Brasil, contacto de emergência e associação a equipas e subequipas.

  \item[Férias e ausências.] Criação e gestão de pedidos com cinco estados explícitos
        (rascunho, pendente, aprovado, rejeitado, cancelado), cálculo de saldo por
        regras contratuais, deteção de sobreposições de ausências na mesma equipa,
        venda de dias, crédito de dias de férias extra e suporte a meio-dias.

  \item[Governação de acessos.] Categorias de permissão com matriz de acesso por função,
        concessão e revogação com justificativa obrigatória, atualizações em lote por
        equipa ou subunidade e visibilidade de permissões por colaborador.

  \item[Auditoria e notificações.] Histórico de versões de pedidos, registo de alterações
        de permissões, notificações de eventos críticos e exportação de dados para
        suporte à decisão.
\end{description}

\subsection{Priorização MoSCoW}

Os requisitos foram classificados com o método MoSCoW para garantir foco nas
funcionalidades de maior valor operacional:

\begin{itemize}
  \item \textbf{Must have:} módulo de perfil de colaborador, fluxo de férias com
        aprovação e histórico, suporte PT/BR nos fluxos principais;
  \item \textbf{Should have:} modelo de permissões com registo de concessão/revogação,
        notificações de eventos críticos;
  \item \textbf{Could have:} relatórios de exportação, banco de horas, dashboards de decisão;
  \item \textbf{Won't have} (no âmbito deste estágio): integração com sistemas de terceiros,
        módulo de \textit{payroll}, aplicação móvel nativa.
\end{itemize}

\subsection{Modelação de Domínio}

A modelação de domínio foi uma das atividades com maior impacto na qualidade final
da solução. A compreensão de que um pedido de férias não é um objeto simples -
mas sim uma entidade que evolui ao longo do tempo, com versões, aprovações associadas
e crédito de saldo - obrigou a um desenho cuidadoso das entidades centrais.

As entidades principais do domínio são:

\begin{itemize}
  \item \textbf{User} e \textbf{Profile} - identidade de acesso e dados de colaborador,
        separados por razões de segurança e de gestão de propriedade de dados;
  \item \textbf{Team} e \textbf{Subteam} - estrutura organizacional que determina
        visibilidade e hierarquia de aprovação;
  \item \textbf{Vacation} - pedido de férias/ausência, com versões que preservam
        o histórico de cada alteração;
  \item \textbf{VacationApproval} - registo de cada etapa de aprovação, com responsável,
        estado e \textit{timestamp};
  \item \textbf{Permission} e \textbf{UserPermission} - modelo de permissões com
        trilho auditável de concessões e revogações;
  \item \textbf{Notification} - alertas operacionais para eventos como mudanças de
        permissão ou resposta a pedidos;
  \item \textbf{VacationBalanceCredit} - crédito de saldo de férias, com suporte a
        crédito manual, venda de dias e dias extra da empresa.
\end{itemize}

A decisão de separar \texttt{User} de \texttt{Profile} foi uma das mais relevantes do
projeto. Permite que dados sensíveis de perfil, como IBAN ou NIF, sejam
protegidos por políticas de \textit{ownership} independentes das credenciais de acesso,
o que é uma exigência real de segurança para informação desta natureza.

\section{Implementação e Desenvolvimento}

O desenvolvimento foi organizado em domínios funcionais, cada um com as suas camadas
de \textit{frontend}, \textit{backend} e persistência. A ordem de implementação seguiu
aproximadamente a prioridade MoSCoW: o módulo de perfil foi o primeiro a ser
desenvolvido, seguido pelo módulo de férias, que foi o mais extenso e o que exigiu
mais iterações, e depois pelo modelo de permissões.

\subsection{Módulo de Perfil de Colaborador}

O módulo de perfil foi o ponto de entrada do projeto. O objetivo era criar um sistema
central de gestão de dados de colaboradores que substituísse as folhas de cálculo
existentes e garantisse validação e segurança desde o início.

A implementação incluiu:

\begin{itemize}
  \item campos de dados pessoais e contratuais com validação explícita;
  \item suporte a campos específicos por país: NIF e NISS para Portugal, CPF,
        PIS e CTPS para o Brasil;
  \item políticas de \textit{ownership} que garantem que dados sensíveis - como
        IBAN ou dados de identificação - só podem ser alterados por utilizadores
        com permissão adequada, e não pelo próprio colaborador sem aprovação;
  \item associação a equipas e subequipas com validação de hierarquia.
\end{itemize}

Uma dificuldade encontrada nesta fase foi a definição clara de quem pode alterar o quê.
A regra geral parece simples: cada um gere o seu perfil. Mas, na prática, há campos
que só o RH pode alterar (dados de contrato), campos que o colaborador pode propor mas
que precisam de confirmação (dados bancários), e campos que qualquer um pode ver mas
ninguém pode alterar diretamente (estrutura de equipa). Este modelo de \textit{ownership}
granular foi implementado no \textit{middleware} de autorização e reutilizado em outros
módulos.

\begin{figure}[H]
  \centering
  % INSTRUÇÃO: Substituir pelo print real da página de perfil do colaborador
  % \includegraphics[width=\textwidth]{imagens/profile_page.png}
  \missingfigure[figwidth=\textwidth]{Interface do módulo de perfil de colaborador}
  \caption{Módulo de perfil de colaborador - vista de detalhe}
  \label{fig:perfil}
\end{figure}

\subsection{Fluxos de Férias e Ausências}
\label{subsec:ferias}

O módulo de férias foi, sem dúvida, o mais complexo e o que consumiu mais tempo de
desenvolvimento. A razão é simples: as regras de negócio para férias não são apenas
complexas em termos de estado: são diferentes para Portugal e para o Brasil, e
têm de coexistir no mesmo sistema.

\subsubsection{Estados e Versionamento}

Cada pedido de férias passa por um ciclo de vida com cinco estados possíveis:

\begin{itemize}
  \item \textbf{Rascunho} - pedido criado mas ainda não submetido;
  \item \textbf{Pendente} - pedido submetido e aguarda aprovação;
  \item \textbf{Aprovado} - pedido aprovado por todos os níveis de aprovação;
  \item \textbf{Rejeitado} - pedido rejeitado por um aprovador;
  \item \textbf{Cancelado} - pedido cancelado pelo próprio colaborador.
\end{itemize}

A decisão de implementar \textbf{versionamento de pedidos} foi tomada cedo no
desenvolvimento e provou ser uma das mais importantes. Em vez de simplesmente
atualizar um registo quando um pedido é alterado, o sistema cria uma nova versão
sempre que há uma mudança significativa. Isto garante um histórico completo e
auditável: é possível saber exatamente como estava o pedido em cada momento, quem
o aprovou em que versão, e que alterações foram feitas após a aprovação.

\subsubsection{Diferenças PT/BR no Fluxo de Aprovação}

A gestão das diferenças legislativas entre Portugal e Brasil foi o aspeto mais
desafiante do módulo de férias. As principais diferenças implementadas foram:

\begin{itemize}
  \item \textbf{Níveis de aprovação.} Em Portugal, um pedido de férias requer aprovação
        de um único nível, o gestor direto. No Brasil, o fluxo é de dois níveis:
        primeiro o gestor de equipa, depois o RH. Esta diferença implicou que a
        entidade \texttt{VacationApproval} tivesse de suportar múltiplos registos
        por pedido, com ordem e responsável distintos consoante o país do colaborador.

  \item \textbf{Prazo de submissão.} Em Portugal, existe uma data-limite de 30 de abril
        para a submissão do mapa de férias. No Brasil, as alterações a pedidos aprovados
        exigem uma antecedência mínima de dez dias úteis.

  \item \textbf{Elegibilidade e saldo.} O cálculo do saldo de férias segue regras
        contratuais distintas. Em Portugal, o trabalhador tem direito a 22 dias úteis
        por ano, com elegibilidade imediata após o período de experiência. No Brasil,
        o cálculo é baseado em meses de empresa e envolve regras específicas do regime CLT.

  \item \textbf{Suporte a meio-dias.} Ambos os contextos suportam pedidos de meio-dia,
        com distinção entre o período da manhã e da tarde.
\end{itemize}

A abordagem adotada para gerir esta complexidade foi criar serviços separados para
as regras de cada país, \texttt{VacationRulesPT} e \texttt{VacationRulesBR},
com uma interface comum que o serviço principal invoca com base no atributo
\texttt{country} do perfil do colaborador. Esta separação torna as regras explícitas,
testáveis de forma independente e fáceis de evoluir sem risco de interferência
entre os dois contextos.

\subsubsection{Validação de Capacidade}

Além das regras por país, o sistema implementa uma validação de capacidade de equipa:
não é possível aprovar um pedido de férias se mais de um terço da equipa estiver de
ausência no mesmo período. Esta regra foi implementada como uma verificação transacional
executada no momento da aprovação, consultando os pedidos aprovados e pendentes
de todos os membros da equipa para o período em questão.

\subsubsection{Cálculo de Saldo}

O saldo de férias de cada colaborador é calculado dinamicamente a partir de quatro fontes:

\begin{itemize}
  \item dias de direito contratual para o ano em curso;
  \item dias utilizados em pedidos aprovados;
  \item dias de crédito extra atribuídos manualmente pela empresa;
  \item dias resultantes de venda de férias ou transferência de saldo.
\end{itemize}

A entidade \texttt{VacationBalanceCredit} foi criada especificamente para registar
as operações sobre o saldo (crédito, débito e venda) de forma auditável e
independente dos pedidos de férias.

\begin{figure}[H]
  \centering
  % INSTRUÇÃO: Substituir pelo print real do módulo de férias (página de listagem ou detalhe)
  % \includegraphics[width=\textwidth]{imagens/vacations_page.png}
  \missingfigure[figwidth=\textwidth]{Interface do módulo de gestão de férias}
  \caption{Módulo de férias - listagem de pedidos com indicador de estado}
  \label{fig:ferias}
\end{figure}

\subsection{Governação de Permissões e Acessos}

O módulo de permissões foi desenhado com dois objetivos simultâneos: flexibilidade
para atribuir diferentes níveis de acesso a diferentes utilizadores, e rastreabilidade
absoluta de todas as alterações.

A implementação inclui:

\begin{itemize}
  \item categorias de permissão organizadas em áreas funcionais (perfil, férias,
        permissões, relatórios);
  \item concessão e revogação de acesso total ou restrito, com justificativa obrigatória
        em cada operação;
  \item atualizações em lote para atribuição de permissões a todos os membros de
        uma equipa de uma só vez;
  \item \textit{engine} de permissões central (\texttt{permission-engine}) que é
        consultado em cada pedido à API para verificar se o utilizador tem acesso
        ao recurso solicitado;
  \item painéis de inspeção que permitem ao RH ver, em tempo real, o conjunto de
        permissões de qualquer colaborador.
\end{itemize}

A decisão de centralizar a lógica de verificação de permissões num único módulo, em vez de replicar verificações em cada rota, foi uma das que mais contribuiu
para a consistência do sistema. Qualquer rota que necessite verificar permissões
invoca o mesmo \textit{engine}, o que garante que as regras são aplicadas de forma
uniforme em toda a aplicação.

\subsection{Camada de Persistência e Integridade de Dados}

A persistência foi implementada com PostgreSQL e Prisma ORM. O Prisma foi escolhido
pela sua integração nativa com TypeScript. Os tipos gerados a partir do schema
são utilizados diretamente no código, o que elimina uma classe inteira de erros de tipo
em tempo de execução, além de facilitar a gestão de migrações.

As principais decisões de implementação na camada de persistência foram:

\begin{itemize}
  \item uso de \textbf{transações} para operações que envolvem múltiplas entidades,
        como a criação de uma nova versão de pedido de férias com atualização do
        saldo associado;
  \item definição de \textit{constraints} e relações no schema que garantem integridade
        referencial mesmo em caso de falha de aplicação;
  \item criação de \textit{scripts} de \textit{seed} para popular o ambiente de
        desenvolvimento com dados representativos, o que foi fundamental para o
        desenvolvimento de testes.
\end{itemize}

\subsection{Desenvolvimento de API e Serviços}

O \textit{backend} foi estruturado segundo o padrão \texttt{route → service → repository}.
As rotas Express recebem os pedidos HTTP, validam a autenticação e autorização,
e delegam ao serviço correspondente. O serviço implementa a lógica de negócio
e delega ao repositório para as operações de base de dados. Esta separação foi
adotada desde o início, mas foi reforçada ao longo do projeto à medida que a
complexidade de cada domínio cresceu.

Os endpoints principais cobrem:

\begin{itemize}
  \item gestão de utilizadores e perfis (\texttt{/users}, \texttt{/profile});
  \item gestão de férias, aprovação e versionamento (\texttt{/vacations});
  \item governação de permissões (\texttt{/permissions});
  \item notificações (\texttt{/notifications});
  \item gestão de ficheiros e \textit{uploads} (\texttt{/files}).
\end{itemize}

A autenticação é baseada em \gls{jwt}: após o \textit{login}, o cliente recebe um
\textit{token} que é incluído em cada pedido subsequente. O \textit{middleware} de
autenticação valida o \textit{token} e extrai o identificador do utilizador, que é
depois usado pelo \textit{engine} de permissões para verificar o acesso ao recurso
solicitado.

\subsection{Interface de Utilizador}

O \textit{frontend} foi desenvolvido em React com TypeScript, com ênfase na clareza
dos fluxos e na redução do esforço cognitivo do utilizador. As páginas com maior
complexidade, como gestão de férias e permissões, foram as que mais iterações
exigiram, porque a apresentação da informação de estado tem impacto direto na
qualidade das decisões dos gestores.

As principais decisões de interface foram:

\begin{itemize}
  \item indicadores de estado visuais e inequívocos em todos os pedidos de férias;
  \item modais de confirmação para ações destrutivas ou irreversíveis, como a
        revogação de uma permissão ou o cancelamento de um pedido aprovado;
  \item \textit{pickers} de colaboradores e equipas com pesquisa e seleção em lote,
        essenciais para as operações de atribuição de permissões;
  \item mensagens de erro e de estado orientadas para o utilizador, não para o
        sistema.
\end{itemize}

\begin{figure}[H]
  \centering
  % INSTRUÇÃO: Substituir pelo print real do módulo de permissões
  % \includegraphics[width=\textwidth]{imagens/permissions_page.png}
  \missingfigure[figwidth=\textwidth]{Interface do módulo de governação de permissões}
  \caption{Módulo de permissões - painel de inspeção de acessos por colaborador}
  \label{fig:permissoes}
\end{figure}

\section{Validação e Controlo de Qualidade}

A estratégia de validação foi definida desde o início: cada módulo devia ter cobertura
de testes antes de ser considerado entregue. Na prática, isto nem sempre foi possível
com a profundidade desejada. A complexidade das regras PT/BR de férias, em particular,
exigiu múltiplas iterações de implementação antes de ser possível escrever testes
suficientemente estáveis.

\subsection{Testes Unitários}

Foram escritos testes unitários para as funções de maior criticidade de negócio,
com foco nas áreas onde um erro tem consequências diretas no utilizador:

\begin{itemize}
  \item funções de cálculo de saldo de férias, para PT e BR;
  \item validações de campo no módulo de perfil (formato de NIF, IBAN, CPF);
  \item lógica de deteção de sobreposição de ausências;
  \item componentes React críticos, com verificação de estados de erro e de carregamento.
\end{itemize}

\subsection{Testes de Integração}

Os testes de integração focaram-se na validação dos fluxos completos entre a API
e a base de dados, incluindo:

\begin{itemize}
  \item fluxo de criação, aprovação e cancelamento de um pedido de férias;
  \item atualização de perfil com validação de regras de \textit{ownership};
  \item atribuição e revogação de permissões com verificação da persistência
        da justificativa.
\end{itemize}

Estes testes foram particularmente importantes para validar o comportamento
das transações da base de dados em cenários de falha parcial.

\subsection{Testes \textit{End-to-End}}

Os testes E2E, implementados com Playwright, cobriram os fluxos de utilizador
mais críticos do ponto de vista operacional:

\begin{itemize}
  \item submissão de um pedido de férias e aprovação por gestor;
  \item criação de perfil de colaborador e visualização de histórico;
  \item alteração de permissões e verificação do efeito na visibilidade de recursos.
\end{itemize}

\subsection{Revisões e Qualidade de Código}

Além dos testes automáticos, foi adotada uma prática de revisão de código antes de
cada integração de funcionalidade, verificando aderência às convenções de nomenclatura
e organização de módulos, separação de responsabilidades entre camadas e ausência
de lógica de negócio diretamente em rotas Express.

\section{Contratempos e Lições Aprendidas}

Três dificuldades foram particularmente relevantes ao longo do projeto:

\begin{enumerate}
  \item \textbf{Complexidade das regras PT/BR.} As diferenças entre os dois contextos
        legislativos só ficaram completamente claras após várias semanas de implementação.
        A regra da antecedência mínima de dez dias úteis no Brasil, por exemplo, só
        surgiu como requisito após a implementação inicial do fluxo de alteração de
        pedidos. A resposta foi criar uma documentação das regras de negócio:
        o ficheiro \texttt{REGRAS\_PERFIS\_E\_FERIAS.tex}, que passou a ser a
        referência de verdade para o desenvolvimento e os testes.

  \item \textbf{Modelação de dados para versionamento.} A decisão de adicionar
        versionamento de pedidos de férias foi tomada a meio do desenvolvimento
        do módulo, após perceber que a abordagem inicial de atualizar diretamente
        o registo não suportava o histórico necessário para auditoria. A migração
        do modelo de dados existente para o novo esquema foi o momento mais tenso
        do projeto, mas o resultado justificou o esforço.

  \item \textbf{Configuração do ambiente de testes.} A integração entre o Prisma
        e o Vitest para testes de integração com base de dados real exigiu uma
        configuração não-trivial do ambiente. O problema principal era garantir
        que cada teste corria com um estado de base de dados limpo, sem interferência
        entre testes. A solução foi criar uma base de dados de teste isolada e um
        \textit{script} de \textit{setup} que limpa e repopula o estado antes de
        cada suite de testes.
\end{enumerate}

A lição mais importante foi a de que as regras de negócio precisam ser documentadas
explicitamente antes de serem implementadas, não o contrário. Tentar perceber uma
regra a partir do código que a implementa é possível, mas ineficiente. A criação
do ficheiro de regras de negócio no meio do projeto deveria ter sido feita no início.

\section{Documentação e Continuidade}

Com vista à continuidade do projeto após o estágio, foram produzidos os seguintes
artefactos de documentação:

\begin{itemize}
  \item ficheiro de regras de negócio PT/BR (\texttt{REGRAS\_PERFIS\_E\_FERIAS.tex})
        com a especificação detalhada de todas as regras implementadas;
  \item documentação das rotas da API com exemplos de pedido e resposta;
  \item \textit{scripts} de migração e \textit{seed} para replicação do ambiente
        de desenvolvimento;
  \item identificação explícita de pontos de dívida técnica e direções de evolução futura.
\end{itemize}

% =========================================================================
\chapter{Arquitetura da Solução e Validação Técnica}
\label{cap:arquitetura}
% =========================================================================

Este capítulo descreve as decisões de arquitetura do SMARTER HUB e apresenta
as evidências de qualidade técnica da solução. A inclusão de um capítulo dedicado
à arquitetura é justificada pela necessidade de demonstrar que as escolhas técnicas
não foram acidentais, mas sim o resultado de um processo de análise que pesou
alternativas, identificou \textit{trade-offs} e tomou decisões conscientes.

\section{Visão Arquitetural}

O SMARTER HUB é uma aplicação \textit{full-stack} composta por três camadas principais,
representadas na Figura~\ref{fig:arquitetura}:

\begin{figure}[H]
  \centering
  % INSTRUÇÃO: Substituir pelo diagrama real de arquitetura (pode ser exportado do Lucidchart ou criado em TikZ)
  % \includegraphics[width=0.85\textwidth]{imagens/arquitetura_logica.png}
  \missingfigure[figwidth=0.85\textwidth]{Diagrama de arquitetura lógica do SMARTER HUB}
  \caption{Arquitetura lógica do SMARTER HUB}
  \label{fig:arquitetura}
\end{figure}

\begin{description}
  \item[Apresentação.] Aplicação React com TypeScript, construída com Vite e
        implantada no servidor interno da empresa. Consome a API REST e apresenta os fluxos de perfil,
        férias, permissões e notificações.

  \item[Serviço.] API REST implementada com Node.js e Express. Recebe pedidos HTTP,
        valida autenticação e autorização via \gls{jwt}, e orquestra a lógica de
        negócio através de serviços organizados por domínio.

  \item[Persistência.] PostgreSQL com Prisma ORM. Armazena colaboradores, pedidos
        de férias, aprovações, permissões, notificações e ficheiros. Implantada
        no Render.
\end{description}

A separação em três camadas, com responsabilidades claras, foi uma decisão deliberada
desde o início do projeto. Numa aplicação com esta complexidade de regras de negócio,
misturar lógica de domínio com validação HTTP ou com acesso a dados seria uma fonte
garantida de \textit{bugs} difíceis de localizar.

\subsection{Componentes Principais}

A Tabela~\ref{tab:componentes} lista os módulos centrais do sistema e a sua localização
no código-fonte.

\begin{table}[H]
  \caption{Componentes principais do SMARTER HUB}
  \label{tab:componentes}
  \small\sffamily
  \begin{tabularx}{\textwidth}{|l|X|}
    \hline
    \textbf{Componente} & \textbf{Localização / Responsabilidade} \\
    \hline
    Autenticação e sessão & \texttt{backend/src/routes/auth.ts}, \texttt{backend/src/middleware/auth.ts}, 
    \texttt{src/portal/auth-storage.ts} \\
    \hline
    Perfil e colaboradores & \texttt{backend/src/routes/users.ts}, \texttt{backend/src/routes/profile.ts}, \texttt{src/pages/CollaboratorsPage.tsx} \\
    \hline
    Férias e ausências & \texttt{backend/src/routes/vacations.ts}, \newline \texttt{backend/src/repositories/vacations.repository.ts}, \newline \texttt{backend/src/services/vacations/*.ts} \\
    \hline
    Motor de permissões & \texttt{backend/src/lib/permission-engine.js}, \newline \texttt{backend/src/routes/permissions.ts} \\
    \hline
    Notificações e ficheiros & \texttt{backend/src/routes/notifications.ts}, \newline \texttt{backend/src/services/files/file-access.service.ts} \\
    \hline
  \end{tabularx}
\end{table}

\subsection{Subsistemas Transversais}

Para além das três camadas principais, o sistema integra subsistemas transversais:

\begin{itemize}
  \item \textbf{Autenticação e autorização}: \gls{jwt} com validação de \textit{scope}
        e \textit{ownership} em cada pedido à API;
  \item \textbf{Notificações por email}: envio via Microsoft Graph API para eventos
        críticos como aprovação ou rejeição de pedidos e alterações de permissões;
  \item \textbf{Gestão de ficheiros}: \textit{upload} e acesso protegido a documentos
        de suporte, com controlo de acesso por permissão.
\end{itemize}

\section{Modelo de Dados e Governança}

O modelo relacional do SMARTER HUB foi construído com foco em consistência e
auditabilidade. A Figura~\ref{fig:der} apresenta o diagrama entidade-relação
das entidades nucleares.

\begin{figure}[H]
  \centering
  % INSTRUÇÃO: Substituir pelo DER real (exportar do Prisma Studio, DBeaver ou Lucidchart)
  % \includegraphics[width=\textwidth]{imagens/modelo_dados.png}
  \missingfigure[figwidth=\textwidth]{Diagrama entidade-relação do modelo de dados}
  \caption{Modelo relacional do SMARTER HUB - entidades nucleares}
  \label{fig:der}
\end{figure}

A governança de dados foi reforçada com:

\begin{itemize}
  \item \textit{foreign keys} e \textit{constraints} definidas no schema Prisma, garantindo integridade referencial a nível de base de dados;
  \item normalização das entidades para evitar duplicação de dados e inconsistências;
  \item transações para operações que alteram múltiplas entidades relacionadas,
        como a criação de uma versão de pedido de férias com atualização do saldo
        (\texttt{createVacation-\\VersionTransaction}) ou o registo de crédito de
        saldo com notificação associada
        (\texttt{createVacationBalanceCreditsWithNotifications});
  \item auditoria implícita através do versionamento de pedidos e do registo
        de todas as alterações de permissões.
\end{itemize}

\section{Decisões de Arquitetura e \textit{Trade-offs}}

\subsection{Padrão \texttt{Route → Service → Repository}}

A separação em três camadas no \textit{backend} (rotas, serviços e repositórios) foi adotada desde o início e reforçada ao longo do projeto. A rota trata exclusivamente
da validação HTTP e da autorização; o serviço implementa a lógica de negócio; o
repositório é o único ponto de acesso à base de dados.

O \textit{trade-off} desta abordagem é a necessidade de criar mais ficheiros e mais
\textit{boilerplate}. O benefício é que qualquer função de negócio pode ser testada
isoladamente, sem dependência de uma rota HTTP ou de uma instância real de base de dados.

\subsection{Versionamento vs. Atualização Direta}

A decisão de versionar os pedidos de férias em vez de os atualizar diretamente
foi a que mais impacto teve no modelo de dados. Versionar é mais caro em termos
de armazenamento e de complexidade de consulta, mas é o único mecanismo que
garante auditabilidade completa: sem versões, não é possível saber como estava
um pedido no momento em que foi aprovado, o que inviabiliza qualquer processo
de auditoria sério.

\subsection{Separação \texttt{User} / \texttt{Profile}}

A separação da identidade de acesso (\texttt{User}) dos dados de colaborador
(\texttt{Profile}) aumenta a complexidade das consultas mas permite um controlo
de acesso granular. Em particular, permite que as políticas de \textit{ownership}
sobre dados sensíveis, como dados bancários ou de identificação, sejam
independentes das políticas de autenticação.

\section{Riscos Residuais}

A solução entregue é funcional e cobre os requisitos principais do projeto, mas
mantém riscos técnicos que devem ser endereçados nas fases seguintes:

\begin{itemize}
  \item \textbf{Lógica de negócio em rotas.} Algumas rotas do \textit{backend}
        ainda contêm lógica de negócio que deveria estar nos serviços correspondentes.
        Esta dívida técnica resulta do crescimento orgânico do sistema e deve ser
        endereçada antes da entrada em produção.
  \item \textbf{Páginas React monolíticas.} Algumas páginas do \textit{frontend}
        cresceram para dimensões que dificultam a manutenção. A decomposição em
        componentes e \textit{hooks} reutilizáveis é o próximo passo natural.
  \item \textbf{Cobertura de testes de integração.} Os testes de integração com
        base de dados real precisam de ser expandidos para cobrir mais cenários
        de falha e condições de corrida, especialmente no módulo de férias.
  \item \textbf{Revisão de autorização em endpoints sensíveis.} A validação de
        \textit{ownership} deve ser verificada em todos os endpoints que acedem
        a dados de outros utilizadores.
\end{itemize}

\section{Síntese do Capítulo}

A arquitetura do SMARTER HUB reflete escolhas deliberadas: a separação em camadas,
o versionamento de pedidos, a separação \texttt{User}/\texttt{Profile} e a centralização
da lógica de permissões são todas decisões que aumentam a complexidade a curto prazo
mas reduzem o risco operacional a longo prazo. Os riscos residuais identificados são
reais, mas são o resultado de decisões conscientes sobre o que priorizar dentro do
tempo disponível, não de desconhecimento das alternativas.

% =========================================================================
\chapter{Conclusão}
\label{cap:conclusao}
% =========================================================================

\section{Resumo e Principais Conclusões}

O SMARTER HUB foi desenvolvido para resolver um problema real: a fragmentação dos
processos de gestão de pessoas numa empresa com colaboradores em dois países com
legislações laborais distintas. No final do período de estágio, entregou-se uma plataforma funcional, com
validação técnica por testes e revisão de código, e está em fase de pré-produção.

O trabalho mais difícil não foi a implementação em si, mas a compreensão das regras de
negócio. As diferenças entre o modelo de aprovação português e o brasileiro, as
políticas de \textit{ownership} sobre dados sensíveis e o mecanismo de versionamento
de pedidos são problemas de domínio que exigiram análise cuidadosa antes de qualquer
linha de código. Essa análise e a documentação que dela resultou é provavelmente
o artefacto mais valioso do projeto.

\section{Objetivos Realizados}

\begin{itemize}
  \item \textbf{Centralizar a gestão de colaboradores:} cumprido. Módulo de perfil
        com campos PT/BR, validações e políticas de \textit{ownership} implementados.
  \item \textbf{Implementar fluxos de férias/ausências:} cumprido. Cinco estados
        explícitos, versionamento, validações de capacidade de equipa e suporte
        a regras PT/BR.
  \item \textbf{Suportar dois contextos legislativos:} cumprido. Fluxos de aprovação,
        cálculo de saldo e regras de elegibilidade implementados de forma independente
        para PT e BR.
  \item \textbf{Estruturar modelo de permissões:} cumprido. Motor de permissões
        centralizado com trilho auditável de concessões e revogações.
  \item \textbf{Disponibilizar mecanismos de reporte:} parcialmente cumprido. Notificações e suporte a exportação entregues; dashboards operacionais
        ficam como trabalho futuro.
  \item \textbf{Garantir qualidade técnica:} parcialmente cumprido. Testes unitários e de integração implementados; cobertura E2E com base de dados real precisa ser expandida.
\end{itemize}

\section{Outros Trabalhos Realizados}

Para além dos módulos principais, foram também realizados trabalhos de suporte que
tiveram impacto na qualidade geral da plataforma:

\begin{itemize}
  \item reforço da segurança de autenticação e de controlo de acesso a ficheiros;
  \item melhoria da gestão de sessão no \textit{frontend};
  \item extração inicial de componentes React para reduzir duplicação nas páginas
        mais complexas;
  \item documentação das regras de negócio PT/BR para uso futuro da equipa técnica.
\end{itemize}

\section{Limitações e Trabalho Futuro}

As limitações mais relevantes a endereçar nas fases seguintes são:

\begin{itemize}
  \item completar a migração do \textit{backend} para o padrão
        \texttt{route → service → repository} de forma consistente;
  \item decompor as páginas React de maior dimensão em componentes reutilizáveis;
  \item expandir os testes de integração com base de dados real para cobrir cenários
        de falha e condições de corrida;
  \item completar a validação de autorização por \textit{ownership} em todos os
        \textit{endpoints} sensíveis;
  \item implementar dashboards operacionais e relatórios de mapa de férias automático.
\end{itemize}

\section{Apreciação Final}

Construir uma plataforma \textit{web} de raiz, de forma individual, para uma empresa
real, num período de três meses, é uma experiência diferente de qualquer projeto
académico. A diferença não está na dificuldade técnica, está na responsabilidade.
As decisões de arquitetura têm consequências reais. Os \textit{bugs} afetam utilizadores
reais. As regras de negócio têm de ser corretas, não aproximadas.

O que mais me surpreendeu foi perceber que as decisões mais difíceis não foram técnicas.
Foram decisões de produto: o que implementar primeiro, o que simplificar, o que adiar.
Gerir essa ambiguidade com informação incompleta, e sem a rede de segurança de um
enunciado com todos os requisitos definidos, foi a aprendizagem mais significativa
deste estágio.

O SMARTER HUB não está terminado. Mas está funcional, está em pré-produção, e tem
uma direção clara de evolução. Para um projeto de três meses desenvolvido de forma
individual, esse é o resultado certo.

% =========================================================================
% BIBLIOGRAFIA
% =========================================================================
\fancyhead[LO,RE]{Bibliografia}
\nocite{*}

% =========================================================================
% ANEXOS
% =========================================================================
\begin{appendices}

\chapter{Regras de Negócio PT/BR - Sumário}
\label{anx:regras}
\fancyhead[LO,RE]{Regras de Negócio PT/BR}

Esta secção apresenta um sumário das principais diferenças de regras de negócio
entre os contextos português e brasileiro, tal como implementadas no SMARTER HUB.
O documento completo de especificação de regras (\texttt{REGRAS\_PERFIS\_E\_FERIAS.tex})
encontra-se no repositório do projeto.

\begin{table}[H]
  \caption{Comparativo de regras de férias PT vs. BR}
  \label{tab:ptbr}
  \small\sffamily
  \begin{tabularx}{\textwidth}{|l|X|X|}
    \hline
    \textbf{Dimensão} & \textbf{Portugal} & \textbf{Brasil} \\
    \hline
    Dias de direito & 22 dias úteis por ano & Calculado por meses de empresa (CLT) \\
    \hline
    Níveis de aprovação & 1 (gestor direto) & 2 (gestor + RH) \\
    \hline
    Prazo de submissão & Mapa anual até 30 de abril & Antecedência mínima de 10 dias úteis para alterações \\
    \hline
    Meio-dias & Suportado (manhã / tarde) & Suportado (manhã / tarde) \\
    \hline
    Venda de dias & Suportada & Suportada com regras específicas CLT \\
    \hline
    Campos de identificação & NIF, NISS & CPF, PIS, CTPS \\
    \hline
  \end{tabularx}
\end{table}

\end{appendices}

\end{document}